% Typisierung ohne Identifikation verschachtelter und flacher Tupel.
\Needspace{14\baselineskip}
\FormulaThmDeltaK[Tafelkandidaten sind verschachtelte natürliche Tupel]{%
 \FiniteSemigroupAlg{A}{\star}\vdash
 n_A\in\mathbb N\land W_A\subseteq V_{n_A}
}{B37TLexCodeCandidatesTyped}{%
 \DeltaRow{Mengen}{A\dsep n_A\dsep W_A\dsep w\dsep i\dsep j}
 \DeltaRow{Funktionen}{e\dsep\star\dsep m\dsep\ell_i}
}
Hier sind \(n_A=\FiniteCard(A)\), \(I=I_{n_A}\),
\(U=T_{n_A}(\mathbb N)\) und \(V_{n_A}=T_{n_A}(U)\).
Im Zweig mit der Nummerierung \(e\) kürzen wir
\(m=m_{\star,e}\) und \(\ell_i=(m(i,j))_{j\in I}\) ab.
Die innere Folge und danach die äußere Folge werden als die durch ihre
Vorschriften eindeutig bestimmten Funktionen gelesen. In beiden Fällen
wird zuerst jeder Wert typisiert; anschließend rechtfertigt
\FormulaRefAuto{TypedFunctionDefinitionPrinciple}[thm] die in
\FormulaRefAuto{FiniteSemigroupTableTupleDef}[def] verwendete Funktion.
Der dort definierte verschachtelte Kandidat bleibt dabei unverändert.
\begin{tabproofwide}
 \proofstepwidestar[1]{\FiniteSemigroupAlg{A}{\star}}{\rA}
 \proofstepwidestar[1]{\Finite A}{\FormulaRefAuto{FiniteSemigroupFinitenessAxiom}[ax]{1}}
 \proofstepwidestar[1]{n_A\in\mathbb N}{\FormulaRefAuto{FiniteCardNatural}[thm]{2}}
 \proofstepwidestar[1]{I\subseteq\mathbb N}{\FormulaRefAuto{PeanoNatSegSubsetN}[thm]{3}}
 \proofstepwidestar[5]{w\in W_A}{\rA}
 \proofstepwidestar[1,5]{\exists e\,(e\colon I\bij A\land w=w_{n_A}(m_{\star,e}))}{\FormulaRefAuto{FiniteSemigroupCanonicalTableCodeDef}[def]{1,5}}
 \proofstepwidestar[7]{e\colon I\bij A\land w=w_{n_A}(m_{\star,e})}{\rA}
 \proofstepwidestar[1,7]{m\colon I\times I\to I}{\FormulaRefAuto{FiniteSemigroupMultiplicationTableDef}[def]{1,3,\rAEa{7}}}
 \proofstepwidestar[9]{i\in I}{\rA}
 \proofstepwidestar[10]{j\in I}{\rA}
 \proofstepwidestar[9,10]{(i,j)\in I\times I}{\FormulaRefAuto{a\in A\dsep b\in B\vdash(a,b)\in A\times B}[thm]{9,10}}
 \proofstepwidestar[1,7,9,10]{m(i,j)\in I}{\FormulaRefAuto{F\colon A\to B\dsep x\in A\vdash F(x)\in B}[thm]{8,11}}
 \proofstepwidestar[1,7,9,10]{m(i,j)\in\mathbb N}{\FormulaRefAuto{A\subseteq B\dsep x\in A\vdash x\in B}[thm]{4,12}}
 \proofstepwidestar[1,7,9]{\forall j\in I\,(m(i,j)\in\mathbb N)}{\rUI{\rRI{10,13}}}
 \proofstepwidestar[1,7,9]{\exists!F\,(F\colon I\to\mathbb N\land\forall j\in I\,(F(j)=m(i,j)))}{\FormulaRefAuto{TypedFunctionDefinitionPrinciple}[thm]{14}}
 \proofstepwidestar[1,7,9]{\ell_i\colon I\to\mathbb N}{\FormulaRefAuto{FiniteSemigroupTableTupleDef}[def]{3,8,9,15}}
 \proofstepwidestar[1,7,9]{\ell_i\subseteq I\times\mathbb N}{\FormulaRefAuto{FunctionGraphSubsetProduct}[thm]{16}}
 \proofstepwidestar[1,7,9]{\ell_i\in\powerset(I\times\mathbb N)}{\FormulaRefAuto{x\in\powerset(A)\eqvdash x\subseteq A}[thm]{17}}
 \proofstepwidestar[1,7,9]{\ell_i\in U}{\FormulaRefAuto{B37TLexTupleDef}[def]{18,16}}
 \proofstepwidestar[1,7]{\forall i\in I\,(\ell_i\in U)}{\rUI{\rRI{9,19}}}
 \proofstepwidestar[1,7]{\exists!G\,(G\colon I\to U\land\forall i\in I\,(G(i)=\ell_i))}{\FormulaRefAuto{TypedFunctionDefinitionPrinciple}[thm]{20}}
 \proofstepwidestar[1,7]{w_{n_A}(m)\colon I\to U}{\FormulaRefAuto{FiniteSemigroupTableTupleDef}[def]{3,8,21}}
 \proofstepwidestar[1,7]{w\colon I\to U}{\rIE{\FormulaRefAuto{a=b\vdash b=a}[thm]{\rAEb{7}},22}}
 \proofstepwidestar[1,7]{w\subseteq I\times U}{\FormulaRefAuto{FunctionGraphSubsetProduct}[thm]{23}}
 \proofstepwidestar[1,7]{w\in\powerset(I\times U)}{\FormulaRefAuto{x\in\powerset(A)\eqvdash x\subseteq A}[thm]{24}}
 \proofstepwidestar[1,7]{w\in V_{n_A}}{\FormulaRefAuto{B37TLexTupleDef}[def]{25,23}}
 \proofstepwidestar[1,5]{w\in V_{n_A}}{\rEE{6,7,26}}
 \proofstepwidestar[1]{W_A\subseteq V_{n_A}}{\FormulaRefAuto{A\subseteq B\coloneq\forall x\,(x\in A\to x\in B)}[def]{\rUI{\rRI{5,27}}}}
 \proofstepwidestar[1]{n_A\in\mathbb N\land W_A\subseteq V_{n_A}}{\rAI{3,28}}
\end{tabproofwide}
Bei \(n_A=0\) sind beide Allquantoren über \(I=\varnothing\) leer.
Das Funktionsdefinitionsprinzip liefert dann die leere äußere Funktion;
keine positive Tupellänge wurde vorausgesetzt.